% PAGES: 161-162
% Continues the Theorem 5.6 proof opened in sec05_c1_3.tex (page 160 / folio
% 144); the line below picks up immediately after that file's equation
% (5.22). No new "\begin{proof}" is opened here.
%
% This file ends mid-derivation, inside a manual "\textit{Solution:}" block
% (the 2x2 spd/spsd characterization example) that is NOT closed on page 162
% -- the source page ends blank after equation (5.27) with no closing
% square, and the derivation continues on PDF page 163. No \begin{...}
% environment is left open (the "Example:"/"Solution:" labels are plain
% manual text, not the \examplex/\answerx environments, matching the
% convention used elsewhere in this book when the label word is not
% "Answer:"), so scripts/check_env_balance.py has nothing to flag here --
% but the next chunk agent must NOT print a fresh "\textit{Solution:}"
% label; it should continue the sentence right after (5.27) and supply the
% closing "\hfill$\square$\par\medskip" once the source's square is found.

where $Q$ is an orthogonal matrix; cf. Theorem 5.4. Let $x \in \R^n$. From (5.22), it
follows that
\begin{equation}
x\tp Ax \;=\; x\tp Q\Lambda Q\tp x \;=\; (Q\tp x)\tp \Lambda (Q\tp x) \;=\; y\tp
\Lambda y,
\end{equation}
where $y = Q\tp x$. If $y = (y_i)_{i=1:n}$, we find from (5.23) that
\begin{equation}
x\tp Ax \;=\; y\tp \Lambda y \;=\; \sum_{i=1}^{n} \lambda_i y_i^2.
\end{equation}

We now prove each part of the theorem by double implication.

(i) ``If $A$ is a symmetric positive semidefinite matrix, then all the eigenvalues of
$A$ are greater than or equal to zero.''

Assume that there exists an eigenvalue $\lambda_j$ of $A$ such that $\lambda_j < 0$,
and let $w_j \neq 0$ be an eigenvector corresponding to $\lambda_j$. Then,
$Aw_j = \lambda_jw_j$ and
\[
w_j\tp Aw_j \;=\; w_j\tp(\lambda_jw_j) \;=\; \lambda_j \cdot w_j\tp w_j \;=\;
\lambda_j\|w_j\|^2 \;<\; 0,
\]
since $\lambda_j < 0$ and $w_j \neq 0$. This contradicts the property (5.19) of a
positive semidefinite matrix.

We conclude that, if $A$ is symmetric positive semidefinite, then all its eigenvalues
must be greater than or equal to zero.

(i) ``If all the eigenvalues of $A$ are greater than or equal to zero, then $A$ is a
symmetric positive semidefinite matrix.''

Assume that $\lambda_i \geq 0$, for all $i = 1:n$. Let $x \in \R^n$ and let
$y = Q\tp x$. If $y = (y_i)_{i=1:n}$, it follows from (5.24) that
\begin{equation}
x\tp Ax \;=\; y\tp \Lambda y \;=\; \sum_{i=1}^{n} \lambda_i y_i^2.
\end{equation}
Since $\lambda_i \geq 0$ for all $i = 1:n$, it follows from (5.25) that $x\tp Ax \geq
0$ for any $x \in \R^n$, and, from Definition 5.7, we conclude that $A$ is a symmetric
positive semidefinite matrix.

(ii) ``If $A$ is a symmetric positive definite matrix, then all the eigenvalues of $A$
are strictly greater than zero.''

Assume that there exists an eigenvalue $\lambda_j$ of $A$ such that $\lambda_j \leq
0$, and let $w_j \neq 0$ be an eigenvector corresponding to $\lambda_j$. Then,
$Aw_j = \lambda_jw_j$ and
\[
w_j\tp Aw_j \;=\; w_j\tp(\lambda_jw_j) \;=\; \lambda_j \cdot w_j\tp w_j \;=\;
\lambda_j\|w_j\|^2 \;\leq\; 0,
\]
since $\lambda_j \leq 0$. This contradicts the property (5.17) of a positive definite
matrix, since $w_j \neq 0$.

We conclude that, if the matrix $A$ is symmetric positive definite, then all its
eigenvalues must be strictly greater than zero.

(ii) ``If all the eigenvalues of $A$ are strictly greater than zero, then $A$ is a
symmetric positive definite matrix.''

Assume that $\lambda_i > 0$, for all $i = 1:n$. Let $x \in \R^n$ and let $y = Q\tp x$.
If $y = (y_i)_{i=1:n}$, it follows from (5.24) that
\begin{equation}
x\tp Ax \;=\; y\tp \Lambda y \;=\; \sum_{i=1}^{n} \lambda_i y_i^2 \;\geq\; 0.
\end{equation}
Since $\lambda_i > 0$ for all $i = 1:n$, it follows from (5.26) that $x\tp Ax = 0$ if
and only if $y_i = 0$ for all $i = 1:n$, i.e., if and only if $y = 0$. Recall that
$y = Q\tp x$, where $Q$ is an orthogonal matrix. Thus,
\[
x\tp Ax = 0 \iff y = 0 \iff Q\tp x = 0 \iff Q(Q\tp x) = 0 \iff x = 0,
\]
since $Q$ is a nonsingular matrix and $QQ\tp = I$; see (5.11).

We conclude that $x\tp Ax > 0$ for any $x \neq 0$, and, from (5.17) of Definition 5.6,
it follows that $A$ is a symmetric positive definite matrix.
\end{proof}

The following result is a simple consequence of Theorem 5.6:

\begin{lemma}
Any symmetric positive definite matrix is nonsingular.
\end{lemma}

\begin{proof}
Let $A$ be a symmetric positive definite matrix. From Theorem 5.6, we obtain that all
the eigenvalues of $A$ are strictly greater than zero. Thus, $0$ is not an eigenvalue
of the matrix $A$, and, from Theorem 4.3, we conclude that the matrix $A$ is
nonsingular.
\end{proof}

\begin{lemma}
The inverse of a symmetric positive definite matrix is also symmetric positive
definite.
\end{lemma}

\begin{proof}
Let $A$ be a symmetric positive definite matrix. Recall from Lemma 5.3 that the
matrix $A$ is nonsingular, and let $A^{-1}$ be the inverse matrix of $A$. From
Lemma 4.3, it follows that, if $\mu$ is an eigenvalue of $A$, then $\frac{1}{\mu}$ is
an eigenvalue of $A^{-1}$. Note that all the eigenvalues of $A$ are strictly greater
than zero, since $A$ is symmetric positive definite; cf. Theorem 5.6. Then, $\mu > 0$,
and therefore $\frac{1}{\mu} > 0$.

In other words, all the eigenvalues of the symmetric matrix $A^{-1}$ are positive,
and, from Theorem 5.6, we conclude that $A^{-1}$ is a symmetric positive definite
matrix.
\end{proof}

In the example below, we find necessary and sufficient conditions for a $2 \times 2$
symmetric matrix to be symmetric positive definite and symmetric positive
semidefinite, by using the eigenvalue related criterion from Theorem 5.6. A proof of
the same results using Sylvester's criterion is given in Lemma 5.5, and is extended to
$3 \times 3$ matrices in Lemma 5.6.

\par\medskip\noindent\textit{Example:}\ (i) The matrix
$\begin{pmatrix} a & b\\ b & d \end{pmatrix}$ is symmetric positive definite if and
only if $a > 0$ and $ad > b^2$.

(ii) The matrix $\begin{pmatrix} a & b\\ b & d \end{pmatrix}$ is symmetric positive
semidefinite if and only if $a \geq 0$, $d \geq 0$, and $ad \geq b^2$.

\par\medskip\noindent\textit{Solution:}\ Let $A = \begin{pmatrix} a & b\\ b & d
\end{pmatrix}$ The characteristic polynomial of the matrix $A$ is
\begin{equation}
P_A(t) \;=\; \det(tI-A) \;=\; t^2 - (a+d)t + ad - b^2.
\end{equation}
